\documentclass[12pt]{article}
\usepackage{amsmath, amssymb}
\usepackage{geometry}
\usepackage{setspace}
\usepackage{graphicx}

\geometry{margin=1in}
\setstretch{1.2}

\title{Geometric Worldline Foundations of the Electron Field}
\author{Jason Duran Dutton}
\date{August 28, 2026}

\begin{document}

\maketitle

\begin{abstract}
The Geometric Worldline Field Model (GWFM) proposes that the electron field of quantum field
theory arises from a single continuous worldline structure embedded in spacetime. This worldline
spans all allowed trajectories, producing a dense presence-density that appears as a continuous field.
Forward-time segments correspond to electrons; backward-time segments correspond to positrons.
Quantum interference results from geometric overlap of worldline segments, while measurement
corresponds to boundary-condition selection of specific segments. GWFM preserves all predictions
of quantum mechanics and QFT while providing a geometric interpretation of superposition, identity,
and field emergence.
\end{abstract}

\section{Introduction}
Quantum mechanics and quantum field theory accurately predict experimental outcomes but leave open the
ontological status of particles, fields, and superposition. GWFM begins with a geometric premise: a continuous
worldline structure spans spacetime, and the electron field is the projection of this structure. Superposition,
interference, and measurement arise from the geometry of this worldline rather than from separate interpretive
rules. This follows the worldline formalism developed by Polyakov \cite{Polyakov1987}.

\section{Core Principles of GWFM}

\subsection{Continuous Worldline Entity}
The electron and positron field arise from a single continuous worldline structure.

\subsection{Time Orientation}
Forward-time segments manifest as electrons; backward-time segments manifest as positrons.

\subsection{Dense Trajectory Ensemble}
The worldline spans all allowed paths, forming a dense geometric structure.

\subsection{Field Emergence}
The electron field is the presence-density of the worldline ensemble.

\subsection{Interference as Overlap}
Interference patterns arise from geometric overlap of worldline segments.

\subsection{Measurement as Boundary Selection}
Measurement restricts the allowed worldline segments.

\subsection{Identity as Global Geometry}
Electron identity is global and geometric, not local or discrete.

\section{Mathematical Framework}

\subsection{Worldline Foundation}
Let spacetime be a manifold $M$ with metric $g_{\mu\nu}$. The entity’s worldline structure is represented by an ensemble
of paths $y(\tau)$ embedded in $M$. A classical relativistic worldline action is


\[
S_{WL} = -m \int d\tau \sqrt{-g_{\mu\nu}(y(\tau)) \dot{y}^\mu(\tau) \dot{y}^\nu(\tau)}.
\]



\subsection{Spin Structure}
Spin degrees of freedom are encoded using Grassmann variables $\psi^\mu(\tau)$:


\[
S_{\text{spin}} = \int d\tau \, \psi_\mu(\tau)\dot{\psi}^\mu(\tau).
\]



\subsection{Electromagnetic Gauge Coupling}
Coupling to an electromagnetic potential $A_\mu(x)$ is


\[
S_{EM} = q \int d\tau \, \dot{y}^\mu(\tau) A_\mu(y(\tau)).
\]



\subsection{Total Worldline Action}


\[
S_{\text{entity}}[y,\psi] = S_{WL}[y] + S_{\text{spin}}[\psi] + S_{EM}[y].
\]



\subsection{Presence Density}


\[
\rho(x) = \int Dy \, D\psi \, e^{iS_{\text{entity}}[y,\psi]} \delta^{(4)}(x - y(\tau)).
\]



\subsection{Emergent Field}


\[
\phi(x) = \int d\tau \, \delta^{(4)}(x - y(\tau)).
\]



\subsection{Curved Spacetime Compatibility}


\[
D_\mu = \partial_\mu + \omega_\mu + iqA_\mu.
\]



\section{Interpretation}

\subsection{One Entity, Many Appearances}
The electron field is the visible projection of one entity’s full path structure.

\subsection{Interference as Path Overlap}
Interference patterns arise where different segments of the worldline overlap in spacetime.

\subsection{Measurement as Path Selection}
Measurement imposes boundary conditions that restrict the allowed paths.

\subsection{Electrons and Positrons}
Forward-time segments appear as electrons; backward-time segments appear as positrons.

\subsection{Spin, Charge, and Fields}
Spin arises from Grassmann variables and geometric rotation. Charge is associated with topological invariants
of the worldline structure.

\section{Superposition and Double-Slit}

\subsection{Superposition as Geometry}
Superposition corresponds to the worldline spanning multiple allowed trajectories.

\subsection{Double-Slit Experiment}
The worldline passes through both slits; the interference pattern is the projection of the worldline ensemble.

\section{Geometric Foundations}

\subsection{Entity as Geometric Object}
The worldline is a geometric set whose projection yields the electron field.

\subsection{Superposition as Geometric Union}


\[
T = T_A \cup T_B.
\]



\subsection{Interference as Intersection}


\[
I(x) = \text{measure}(T_A \cap T_B).
\]



\subsection{Time Reversal as Orientation}
Backward-time segments correspond to


\[
\frac{dt}{d\tau} < 0.
\]



\subsection{Spin as Geometric Rotation}
Spin is rotation in an internal fiber bundle attached to the worldline.

\subsection{Charge as Topological Invariant}
Charge may be associated with winding number or Chern classes.

\subsection{Measurement as Boundary Geometry}
Measurement is the intersection of the worldline with detector boundary geometry.

\subsection{Entanglement as Unified Geometry}
Entangled systems correspond to a single geometric object spanning multiple regions.

\section*{Appendix A: Mathematical Derivations}
Worldline action, path integral measure, presence density, spinor construction, and Dirac field emergence.

\section*{Appendix B: Extended Structures}
Speculative content removed. Physical extensions preserved.

\section*{Appendix C: Superposition – Deep Interpretation}
Superposition as distributed manifestation of one worldline across multiple trajectories.

\begin{thebibliography}{99}

\bibitem{Schubert2001}
C. Schubert, “Perturbative quantum field theory in the string-inspired formalism,”
Phys. Rept. 355, 73–234 (2001).

\bibitem{Strassler1992}
M. J. Strassler, “Field theory without Feynman diagrams: One-loop effective actions,”
Nucl. Phys. B 385, 145–184 (1992).

\bibitem{Polyakov1987}
A. M. Polyakov, “Gauge Fields and Strings,” Harwood Academic Publishers (1987).

\end{thebibliography}

\end{document}

